\section{Introducción}

%\begin{frame}
%\frametitle{Campo de estudio}

%(galaxias, compocision)

%\end{frame}

\begin{frame}
\frametitle{Objetivo}

\note{
%- Galaxias son sistemas estelares complejos con varias componentes estelares que coexisten simultáneamente y que definen su morfología y cinemática. \\
%- Compuestas por núcleo, disco fino y grueso, halo estelar, barra y brazos espirales. \\
\begin{itemize}
    %\item<2-> Resolver problema: Tecnicas de clustering con aprendizaje no supervisado
    %\item<2-> Para estudiar la formación y evolución de las galaxias es indispensable entender el proceso de \textbf{formación} y \textbf{evolución} de cada una de estas \textbf{componentes}.
\end{itemize}

}

Avanzar en la comprensión de las galaxias simuladas desde el punto de vista de sus componentes, utilizando un enfoque multidisciplinario que combine técnicas de machine learning y astronomía.

\pause

\vspace{2cm}

\centering \textbf{Problema: Falta de etiquetas}

\end{frame}

\iflong
\begin{frame}
\frametitle{Aplicaciones}

\begin{itemize}
    \item<1-> Avance en la comprensión de galaxias
    \item<2-> Evaluación de nuevos métodos de clustering
    \item<3-> Metodología de evaluación de resultados
\end{itemize}

\end{frame}
\fi
